\chapter{Conclusions} \label{sec:Conclusions}

This thesis addressed two engineering problems encountered in building the Rapid Urban Forest Assessment (RUFA) dashboard: querying and aggregating millions of tree records across hundreds of California census-designated places in real time, and producing stable spatial summaries at multiple zoom levels without recomputing cluster assignments on every map interaction. Each problem was approached by characterizing the workload, prototyping a straightforward solution, measuring where it broke down, and replacing it with an implementation that scales to the production dataset. The RUFA Score formula itself was developed at UFEI; the contribution here is the infrastructure that computes and serves that score and its maps interactively.

On the data side, RUFA's MySQL schema combines normalized inventory and aerial-detection tables with cascading materialized views that propagate RUFA Score metrics from census places to ZIP codes, tracts, and counties. Spatial indexing strategies were evaluated for assigning tree points to administrative polygons, and indexed queries achieved sub-second response times against the full California inventory. The resulting design preserves compatibility with the existing \href{https://selectree.calpoly.edu}{selectree.calpoly.edu} infrastructure while delivering the interactive performance required by the dashboard.

On the visualization side, the clustering pipeline evolved from a fixed-$k$ K-means baseline, which produced overlapping markers in dense neighborhoods and recomputed assignments on every zoom event, to a SuperCluster index paired with Voronoi tessellation. The final implementation precomputes a multi-resolution cluster hierarchy and derives per-cell tree density and TD-50 diversity scores at query time, yielding zoom-aware summaries that match the uneven density of urban tree distributions.

Together with the AWS-based three-tier deployment described in \autoref{sec:SystemDesign}, these contributions form a system that aggregates inventory and aerial-detection data, computes the four-component RUFA Score, and serves interactive assessments to urban foresters, planners, and researchers. The composite score reduces a complex, multi-source dataset to a single comparable index, supporting evidence-based decisions about planting, maintenance, and policy across California communities.

\section{Future Work}

Several directions extend RUFA beyond the scope of this thesis.

\subsection{Allometric Estimation}

UFEI's AllomeTree tool \cite{ufeiAllometree} exposes species-specific allometric equations for California's urban forest, predicting diameter at breast height (DBH), height, and canopy width with 95\% prediction intervals from either DBH or tree age. The equations derive from the U.S. Forest Service Urban Tree Database \cite{mcpherson2016utd}, which fits 365 equation sets across 171 species and 16 climate regions. Integrating these models into RUFA would let the system estimate structural attributes for trees with incomplete inventory records and for the aerial-only detections produced by the Ventura et al.\ pipeline, which currently provide location and canopy extent but no DBH or height. With those estimates, RUFA could extend its reporting to ecosystem services such as carbon storage, leaf area, and stormwater interception, derived per-tree rather than aggregated from sparse inventory measurements.

\subsection{Incremental Data Integration}

The current materialized-view refresh runs as a batch process. Incremental refresh triggered by inventory updates or new aerial detections would keep RUFA Scores current without recomputing every dependent view, and would support partner organizations contributing inventory data on a rolling basis.

\subsection{Mobile and Field Interfaces}

A field-oriented interface would let urban foresters capture or correct inventory records on site and view RUFA assessments from mobile devices. Bidirectional sync between field captures and the central database would tighten the loop between assessment and on-the-ground management.

\subsection{Geographic Expansion}

Extending RUFA beyond California requires adapting the schema, score normalization, and species references to other climate zones and regulatory contexts. The data model and clustering pipeline are state-agnostic; the main work lies in re-binning each metric against a new statewide distribution and incorporating regional species lists.

As urban areas continue to grow and face increasing environmental pressure, tools that quantify urban forest structure, diversity, and equity will become more important to municipal planning. RUFA contributes to that toolset by combining automated detection, inventory aggregation, and interactive visualization in a single, comparable assessment.
